% Figure for Quantum Mechanics §B2.3 (Example: momentum of the oscillator ground state; Griffiths Problem 3.11): |Phi_0(p)|^2 = exp(-u^2)/sqrt(pi) with u = p/sqrt(m hbar omega), tails |u| > 1 (outside the classical range) shaded, total erfc(1) = 0.157; classical momentum density 1/(pi sqrt(1 - u^2)) of an oscillator with the same energy, dashed. Computed.
\begin{tikzpicture}[font=\small]
  \begin{axis}[nfig, width=9cm, height=5.6cm, xmin=-3.1, xmax=3.1, ymin=0, ymax=1.05,
      xtick={-3,-2,-1,0,1,2,3}, xticklabels={$-3$,$-2$,$-p_{\max}$,$0$,$p_{\max}$,$2$,$3$},
      ytick={0,0.2,0.4,0.6,0.8,1.0},
      xlabel={$p/\sqrt{m\hbar\omega}$}, ylabel={probability density}, clip=true]
    \addplot[draw=none, fill=figR!35, domain=1:3.1, samples=60] {exp(-x^2)/sqrt(pi)} \closedcycle;
    \addplot[draw=none, fill=figR!35, domain=-3.1:-1, samples=60] {exp(-x^2)/sqrt(pi)} \closedcycle;
    \addplot[figO, thick, dashed, domain=-0.9985:0.9985, samples=301] {1/(pi*sqrt(1-x^2))};
    \addplot[figB, very thick, domain=-3.1:3.1, samples=200] {exp(-x^2)/sqrt(pi)};
    \draw[black!55, thin, densely dotted] (axis cs:-1,0) -- (axis cs:-1,1.05);
    \draw[black!55, thin, densely dotted] (axis cs:1,0) -- (axis cs:1,1.05);
    \node[figB, font=\scriptsize, anchor=west] at (axis cs:1.3,0.47) {quantum, $|\Phi_0(p)|^2$};
    \draw[figB, thin] (axis cs:1.3,0.47) -- (axis cs:0.62,0.39);
    \node[figO, font=\scriptsize, anchor=west, align=left] at (axis cs:-2.95,0.86) {classical,\\ same energy};
    \draw[figO, thin] (axis cs:-1.75,0.86) -- (axis cs:-0.94,0.80);
    \node[figR, font=\scriptsize, anchor=south] at (axis cs:2.1,0.16) {$7.9\%$};
    \node[figR, font=\scriptsize, anchor=south] at (axis cs:-2.1,0.16) {$7.9\%$};
    \draw[figR, thin] (axis cs:2.1,0.16) -- (axis cs:1.35,0.05);
    \draw[figR, thin] (axis cs:-2.1,0.16) -- (axis cs:-1.35,0.05);
  \end{axis}
\end{tikzpicture}
